% Etapa 4 do enunciado: projeção em perspectiva.
% Fonte no código: math_3d.py (projecao_perspectiva, divisao_perspectiva), main.py (ndc_para_pixel)
% Figura do item 4a: visualizar_projecao.py (salva pelo botão de salvar da janela).
\subsection{Projeção em perspectiva}\label{sec:projecao}

Na projeção em perspectiva, objetos mais distantes aparecem menores \cite{azevedo2018}. O volume de
visão é definido pelo ângulo de abertura vertical ($fov = 58^\circ$), pela proporção da janela
($1$) e pelas distâncias dos planos próximo ($0{,}1$) e distante ($100$). Dele sai o valor
$f = \cot(fov/2)$, que corresponde ao $d$ da aula: a posição projetada é

\begin{equation}
  x' = d \cdot \frac{x}{z} .
\end{equation}

Como $z$ está no denominador, quanto mais longe o ponto, mais perto do centro ele fica. No código,
essa conta é feita com uma matriz $4 \times 4$ e uma divisão final pela coordenada $w$. O resultado
são coordenadas entre $-1$ e $1$, que depois são convertidas em pixels da janela de projeção.
Triângulos fora do volume de visão ou voltados para trás não são iluminados nem rasterizados
(Seção~\ref{sec:rasterizacao}). A Figura~\ref{fig:projecao} mostra, à esquerda, os sólidos em 3D
no sistema da câmera com a janela de projeção e, à direita, a janela vista de frente.

\figura{03-desenvolvimento/04-projecao/eye1.png}{\textwidth}{Sólidos em 3D e projeção na janela, com $eye = (9{,}75;\ 13;\ 6{,}5)$}{fig:projecao}{Autoria própria (2026)}
